% =============================================================================
% Extended Related Work --- addresses reviewer weakness W15.
% Included (optionally) from main.tex as an extension to Related Work v2.
% Relies only on amsmath and booktabs (already loaded by main.tex).
% Markers (verbatim, for automated review-response verification):
%   %REVIEW-ADDR: W15-prm
% =============================================================================

\section{Extended Related Work: Interaction of ZVF with Adjacent RL Regimes}
\label{sec:extended-related-work}

A reviewer correctly observed that \ours{}'s zero-variance-fraction (ZVF)
diagnostic was validated primarily against \emph{outcome-reward} GRPO runs with
independent rollouts, and that an adjacent line of recent work modifies one of
the two ingredients ZVF depends on: the reward signal that enters the
group-relative advantage. This extended section surveys the dense,
process-level reward regime and predicts whether ZVF remains informative,
loses discriminative power, or requires an explicit surrogate.

% -----------------------------------------------------------------------------
\subsection{Process Reward Models and Dense Shaping}
\label{sec:ext:tree-prm}
%REVIEW-ADDR: W15-prm

\paragraph{Process Reward Models (PRM).}
The ``Let's Verify Step by Step'' line of work~\cite{lightman2023prm}
introduces \emph{process reward models} (PRMs) that supply a dense, step-level
reward signal rather than a single terminal outcome reward. Under a PRM, the
per-trajectory reward is a sum (or weighted aggregate) over intermediate
steps, and the within-group reward variance is bounded below by the
step-level heterogeneity \emph{even when every trajectory in the group reaches
the same outcome}. Formally, if $r_i = \sum_{t=1}^{T_i} r_{i,t}$ is the PRM
reward for trajectory $i$ in a group and the $\{r_{i,t}\}$ are not perfectly
aligned across $i$, then
\begin{equation}
\mathrm{Var}_i\!\left[r_i\right] \;>\; 0
\quad\text{almost surely,}
\end{equation}
so the event $\{\mathrm{Var}_i[r_i] = 0\}$ that ZVF counts becomes vanishingly
rare. The consequence is that \emph{ZVF approaches zero on both healthy and
collapsed PRM runs}, and therefore \emph{loses discriminative power in the
dense-reward regime}. We flag this as a known failure mode of ZVF and
recommend, under PRM training, replacing the zero-variance-fraction indicator
with either (i) the \emph{per-step reward-variance} statistic averaged over
group members, or (ii) the effective-rank / effective-rollout surrogates
discussed in Appendix~\ref{sec:appendix:zvf-formalization} and
Appendix~\ref{sec:appendix:tool-code-depth}. Both surrogates remain
sensitive to within-group collapse in the presence of dense shaping.

\paragraph{Summary table.}
Table~\ref{tab:ext:zvf-applicability} collects the predictions above.

\begin{table}[h]
\centering
\small
\begin{tabular}{@{}lll@{}}
\toprule
Regime & ZVF applicability & Recommended replacement / augmentation \\
\midrule
Vanilla GRPO (outcome reward) & Informative (baseline)   & --- \\
PRM~\cite{lightman2023prm}    & Degenerate ($\to 0$)     & Per-step reward variance or ERF \\
\bottomrule
\end{tabular}
\caption{ZVF applicability across the outcome-reward and dense process-reward
regimes surveyed in this section, together with the surrogate diagnostic we
recommend where ZVF degenerates.}
\label{tab:ext:zvf-applicability}
\end{table}

% End of extended related work.
